\SourceLeaf{523}\par\medskip\Needspace{4\baselineskip}\pdfbookmark[1]{372.}{sectio-octava-art-372}\centerline{\textbf{372.}}\nobreak\smallskip
Ist $n$ eine (ungerade) Primzahl, so bestehen die Reste $m$ aus null, wofür $M=1$ gilt, und $\frac12(n-1)$ anderen Zahlen, für die $M=2$ gilt. Bezeichnen wir diese Reste (ohne den Rest $0$) ohne nähere Festlegung mit $\mu$, so ist unsere Folge $=1+2\Sigma r^\mu$. Bezeichnet ferner $\nu$ ohne nähere Festlegung alle übrigen Zahlen unter $n$, deren Anzahl ebenfalls $\frac12(n-1)$ ist und die alle quadratischen Nichtreste von $n$ unter $n$ umfassen, so gilt offenbar
\[
1+\Sigma r^\mu+\Sigma r^\nu=1+r+rr+r^3\dotsfour+r^{n-1}=\frac{1-r^n}{1-r}=0
\]
Setzen wir also die Summe unserer Folge, das heißt $1+2\Sigma r^\mu=A$, so gilt $1+2\Sigma r^\nu=-A$ und außerdem $\Sigma r^\mu-\Sigma r^\nu=A$.
\par Nach Artikel $356$ gilt daher $A=\pm\surd n$ oder $\pm\surd{-n}$, je nachdem, ob $n$ die Kongruenz $\equiv1$ oder $\equiv3\, (\mathrm{mod}.4)$ erfüllt. Das Vorzeichen der Wurzel ist dadurch aber noch nicht bestimmt.
\par Ersetzt man in unserer Folge, die wir mit $\Pi$ bezeichnen wollen, $r$ durch eine andere gleichartige Wurzel der Gleichung $x^n-1=0$, etwa $r'=r^\lambda$, so sei $\Pi'$ das Ergebnis.
\par\medskip\Needspace{4\baselineskip}\pdfbookmark[1]{373.}{sectio-octava-art-373}\centerline{\textbf{373.}}\nobreak\smallskip
Ist $n$ das Quadrat oder eine höhere Potenz einer Primzahl, etwa $=p^\pi$, so bestehen einige Reste $m$ aus nicht durch $p$ teilbaren Zahlen; andere sind durch $pp$, aber nicht durch eine höhere Potenz von $p$ teilbar; wieder andere lassen sich durch $p^4$, aber nicht. durch $p^5$ teilen, und so weiter bis zu denen, die durch $p^{\pi-2}$, aber nicht durch $p^{\pi-1}$ oder durch $p^{\pi-1}$, aber nicht durch $p^\pi$ teilbar sind, je nachdem, ob $\pi$ gerade oder ungerade ist. Schließlich kommt der Rest $0$ hinzu, der einzige durch $p^\pi$ teilbare Rest (vgl. Artikel $102$). Bezeichnen wir nun ohne nähere Festlegung mit $\mu$ die quadratischen Reste der Zahl $p$ unter $p$ ohne null (deren Anzahl $=\frac12(p-1)$ ist), so lassen sich diese verschiedenen Restklassen wie folgt darstellen. Die ersten, die nicht durch $p$ teilbar sind, werden durch $\mu+kp$ dargestellt; dabei sind für $k$ alle ganzen Zahlen von $0$ bis $p^{\pi-1}-1$ einzusetzen, sodass die Anzahl aller in dieser Form enthaltenen Reste $=\frac12(p-1)p^{\pi-1}$ ist; für jeden von ihnen gilt $M=2$. Die Summe aller diesen Resten entsprechenden Terme in $\Pi$ ist
\[
\begin{aligned}
&=2\Sigma r^{\mu+kp}=2\Sigma r^\mu\pdot\Sigma r^{kp}\\
&=2\Sigma r^\mu\pdot\frac{r^{p^\pi}-1}{r^p-1}=0
\end{aligned}
\]
Die zweite Klasse der Reste wird durch $\mu pp+kp^3$ dargestellt; dabei sind für $k$ alle ganzen Zahlen von $0$ bis $p^{\pi-3}-1$ einzusetzen, sodass die Anzahl aller in dieser
